Stack Allocator (стековый аллокатор) основан на принципе стека: память выделяется и освобождается по принципу LIFO (Last In, First Out), то есть последний выделенный блок памяти должен быть освобождён первым. Фактически он является продвинутой версией линейного аллокатора, в котором можно освободить некоторый блок памяти, выделенный последним, а не сразу всю отданную аллокатору память. В то же время он сохраняет принцип пространственной локальности и обеспечивает минимальную фрагментацию.

При инициализации стековый аллокатор получает большой непрерывный блок памяти. Внутри аллокатора хранится указатель (<<вершина стека>>), который указывает на начало свободной области памяти. Когда программа запрашивает память, аллокатор перемещает указатель вперёд на размер запрашиваемого блока и возвращает адрес начала этого блока.

Память освобождается в порядке, обратном выделению. Чтобы освободить последний выделенный блок, необходимо переместить указатель назад на размер этого блока. Для освобождения произвольного блока памяти придётся освободить все блоки, выделенные после данного.

Применения стекового аллокатора:

\begin{itemize}
	\item В рекурсивных алгоритмах. Здесь сначала происходит спуск вглубь рекурсии, а затем возврат из неё в порядке, обратном спуску.
	\item В графовых алгоритмах, например, в алгоритме поиске пути A*, где сначала создаются временные узлы, а после завершения они освобождаются.
\end{itemize}
